\documentclass[preview]{standalone}

\usepackage{amsmath}
\usepackage{amssymb}
\usepackage{stellar}
\usepackage{definitions}

\begin{document}

\id{james-funzionalismo}
\genpage

\section{James e il funzionalismo}

\begin{snippetdefinition}{funzionalismo-definition}{Funzionalismo}
    Il \emph{funzionalismo} è un orientamento psicologico che attribuisce
    importanza alle funzioni della mente e alle abitudini apprese. Studia i
    processi mentali in rapporto alla loro utilità adattiva, cioè al modo in
    cui permettono all'organismo di adattarsi all'ambiente e di funzionare con
    efficacia.
\end{snippetdefinition}

\begin{snippet}{william-james-funzionalismo}
    Negli Stati Uniti la figura centrale del funzionalismo è
    \emph{William James}. La psicologia funzionale e descrittiva si oppone a
    Wundt e Titchener perché è più interessata alle funzioni della mente che
    alla sua struttura. Il punto di partenza è la teoria dell'evoluzione e il
    legame tra ambiente e individuo.
\end{snippet}

\begin{snippet}{james-idee-principali}
    Secondo James:
    \begin{itemize}
        \item la psicologia è la scienza della vita mentale, dei suoi fenomeni
        e delle sue condizioni;
        \item l'oggetto della psicologia è l'esperienza immediata, non una
        combinazione di elementi individuati tramite introspezione;
        \item le sensazioni sono l'esito di un processo articolato di inferenza
        e astrazione;
        \item la coscienza è un flusso continuo, condizionato dal cervello;
        \item il metodo comparativo riduce la soggettività dell'introspezione;
        \item i fenomeni psichici sono funzioni dell'organismo finalizzate
        all'adattamento.
    \end{itemize}
\end{snippet}

\begin{snippet}{funzionalismo-applicazione}
    Mentre lo strutturalismo difende una psicologia pura e sperimentale, i
    funzionalisti sostengono anche una psicologia applicata. La spiegazione di
    un comportamento deve considerare il contesto in cui il comportamento si
    manifesta. Con il funzionalismo cambia anche il laboratorio: l'oggetto
    d'indagine non è solo la coscienza, ma il comportamento nella sua totalità.
\end{snippet}

\end{document}
